% XeLaTeX: шрифты через fontspec, русский язык через polyglossia.
% Конфигурация по канону Russian-Phd-LaTeX-Dissertation-Template (А. Акиньшин):
% явно объявляется \cyrillicfont для polyglossia, иначе переносы и метрики кириллицы
% подхватываются нестабильно.
\usepackage{fontspec}
\defaultfontfeatures{Ligatures=TeX}

\usepackage{polyglossia}
\setmainlanguage{russian}
\setotherlanguage{english}

\setmainfont{Times New Roman}
\newfontfamily\cyrillicfont{Times New Roman}
\newfontfamily\englishfont{Times New Roman}

% Моноширинный шрифт с поддержкой кириллицы (для \texttt{...} рядом с кириллическим текстом).
\setmonofont{Courier New}
\newfontfamily\cyrillicfonttt{Courier New}

% Математические пакеты (\mathbb, \varnothing, \mathcal, окружения align/equation и т.п.)
\usepackage{amsmath}
\usepackage{amssymb}

% Русская типография: одиночный пробел после точки (вместо English-style увеличенного).
\frenchspacing

% Умеренная защита от overfull \hbox. Жёсткие настройки не нужны:
% переносы корректно работают через polyglossia с русскими паттернами.
\emergencystretch=1em

% Уменьшение вертикальных отступов до и после display-формул:
% по умолчанию у extbook 14pt они слишком велики, а \onehalfspacing их дополнительно растягивает.
\setlength{\abovedisplayskip}{8pt plus 2pt minus 2pt}
\setlength{\belowdisplayskip}{8pt plus 2pt minus 2pt}
\setlength{\abovedisplayshortskip}{4pt plus 2pt minus 2pt}
\setlength{\belowdisplayshortskip}{4pt plus 2pt minus 2pt}

% Настройка полей: левое 30 мм, правое 15 мм (допускается 10 мм), верхнее 20 мм, нижнее 20 мм (п. 4.2)
\usepackage[a4paper,top=20mm,bottom=20mm,left=30mm,right=15mm]{geometry}

% Межстрочный интервал 1.5, размер шрифта 14 пт (п. 4.2)
\usepackage{setspace}
\onehalfspacing

% Абзацный отступ 1.25 см (п. 4.2)
\setlength{\parindent}{1.25cm}
\setlength{\parskip}{0pt}

% Нумерация страниц: арабскими цифрами, по центру нижней части, без точки (п. 4.3.1)
% Нумерация сквозная, включая приложения — см. \setcounter{page} в main.tex
\usepackage{fancyhdr}
\pagestyle{fancy}
\fancyhf{}
\renewcommand{\headrulewidth}{0pt}
\cfoot{\thepage}
\fancypagestyle{plain}{%
  \fancyhf{}%
  \renewcommand{\headrulewidth}{0pt}%
  \cfoot{\thepage}%
}

% Настройка заголовков через titlesec
\usepackage{titlesec}
% Главы: прописными, полужирный, по центру, без точки (п. 4.4.1, 4.4.4)
\titleformat{\chapter}[hang]{\bfseries\fontsize{16pt}{18pt}\selectfont\centering}{\thechapter}{1em}{\MakeUppercase}
% Разделы: полужирный, с абзацного отступа (п. 4.4.4)
\titleformat{\section}[hang]{\bfseries\normalsize}{\thesection}{1em}{}
% Подразделы: обычный шрифт, с абзацного отступа (п. 4.4.4 — полужирный только для разделов)
\titleformat{\subsection}[hang]{\normalfont}{\thesubsection}{1em}{}
\titlespacing{\chapter}{0pt}{0pt}{20pt}
\titlespacing{\section}{1.25cm}{12pt}{6pt}
\titlespacing{\subsection}{1.25cm}{6pt}{3pt}

% Глубина содержания и нумерации
\setcounter{tocdepth}{2}      % в содержании: главы + разделы + подразделы
\setcounter{secnumdepth}{3}   % нумерация до уровня подпунктов (1.1.1)

% Настройка содержания: обычный шрифт (не жирный) (п. 4.2), точечные лидеры (п. 3.4)
\usepackage{tocloft}
\renewcommand{\cftdotsep}{1}
\renewcommand{\cftchapleader}{\cftdotfill{\cftdotsep}}
\renewcommand{\cftchapfont}{\normalfont}
\renewcommand{\cftchappagefont}{\normalfont}
\renewcommand{\cftsecfont}{\normalfont}
\renewcommand{\cftsecpagefont}{\normalfont}
\renewcommand{\cftsubsecfont}{\normalfont}
\renewcommand{\cftsubsecpagefont}{\normalfont}

% Конфигурационные файлы
\usepackage{enumitem}
\setlist{nosep} % Убирает дополнительные отступы между элементами списков

% По ГОСТ 7.32 и требованиям ИТМО (Приложение 1) маркированные списки — через тире.
% Маркер выставляется с абзацного отступа 1.25 см (норма ИТМО п. 4.5).
\setlist[itemize]{leftmargin=1.25cm, label=--}

% Нумерованные списки: арабские цифры со скобкой: 1), 2), 3) ...
\setlist[enumerate]{leftmargin=1.25cm, label=\arabic*)}

\input{config-formulas.tex}
% Подключаем нужные пакеты для таблиц и оформления
\usepackage{booktabs}   % Для красивых горизонтальных линий (\toprule, \midrule, \bottomrule)
\usepackage{multirow}   % Для объединения ячеек по вертикали
\usepackage{caption}    % Для \ContinuedFloat и настройки заголовков
\usepackage{longtable}
\usepackage{etoolbox}

\DeclareCaptionLabelSeparator{emdash}{\space---\space}  % пробел, длинное тире, пробел

% Настройка заголовков таблиц: слева без абзацного отступа, разделитель — тире (п. 4.6.3)
\captionsetup[table]{position=t, singlelinecheck=false, justification=raggedright, labelsep=emdash, name=Таблица}

% Переключатель для смены выравнивания заголовков при переносе таблицы (п. 4.6.3)
\newcommand*\RightCaption{\captionsetup{justification=raggedleft}}
\newcommand*\LeftCaption{\captionsetup{justification=raggedright}}

% ------------------------
% Шаблон: Многоуровневый заголовок
% ------------------------
\newcommand{\ExampleMultiLevelTable}{%
\begin{table}[h!]
\caption{Пример построения таблицы}
\centering
\begin{tabular}{|c|c|c|c|}
\hline
\multirow{2}{*}{\textbf{Заголовок}} & \multicolumn{2}{c|}{\textbf{Заголовок колонки}} & \multirow{2}{*}{\textbf{Заголовок колонки}}\\
\cline{2-3}
& \textbf{Подзаголовок} & \textbf{Подзаголовок} & \\
\hline
A & B & C & D \\
\hline
\end{tabular}
\end{table}
}

% Шаблон: «Разбитая» таблица (две части под одним номером)
% Номер таблицы в строке «Продолжение таблицы N» подставляется автоматически
\newcommand{\ExampleSplitTable}{%
  % --- Первая часть ---
  \begin{table}[h!]
    \LeftCaption
    \caption{Пример переноса таблицы}
    \label{tab:split_example}
    \centering
    \begin{tabular}{|c|c|c|c|}
    \hline
    \textbf{Заголовок} & \textbf{Заголовок колонки} & \textbf{Заголовок колонки} & \textbf{Заголовок колонки} \\
    \hline
    1 & 2 & 3 & 4 \\
    \hline
    2 & 3 & 4 & 5 \\
    \hline
    \end{tabular}
  \end{table}

  % --- Продолжение (номер подставляется автоматически через \thetable) ---
  \begin{table}[h!]
    \ContinuedFloat
    \RightCaption
    \caption*{Продолжение таблицы \thetable}
    \centering
    \begin{tabular}{|c|c|c|c|}
    \hline
    \textbf{Заголовок} & \textbf{Заголовок колонки} & \textbf{Заголовок колонки} & \textbf{Заголовок колонки} \\
    \hline
    3 & 4 & 5 & 6 \\
    \hline
    4 & 5 & 6 & 7 \\
    \hline
    \end{tabular}
  \end{table}
}

\usepackage{graphicx}
\graphicspath{{images/}} % Папка с изображениями
% caption и emdash-разделитель уже объявлены в config-tables.tex

% Настройка подписей к рисункам: по центру, под рисунком, разделитель — тире (п. 4.5.3)
\captionsetup[figure]{position=b, singlelinecheck=false, labelsep=emdash, justification=centering, name=Рисунок}


% Библиография: ГОСТ-numeric, порядок по первому цитированию (п. 4.10.3–4.10.5)
\usepackage[backend=biber,style=gost-numeric,sorting=none]{biblatex}
\addbibresource{bibliography.bib}

% Убираем курсив из названий источников
\DeclareFieldFormat*{title}{#1}
\DeclareFieldFormat*{journaltitle}{#1}
\DeclareFieldFormat*{booktitle}{#1}
\DeclareFieldFormat*{issuetitle}{#1}
\DeclareFieldFormat{namewrapperformat}{#1}

% Убираем курсив из имён авторов (\mkgostheading — внутренняя команда gost-numeric)
\renewcommand*{\mkgostheading}[1]{#1}
\renewcommand*{\mkbibnamefamily}[1]{#1}
\renewcommand*{\mkbibnamegiven}[1]{#1}
\renewcommand*{\mkbibnameprefix}[1]{#1}
\renewcommand*{\mkbibnamesuffix}[1]{#1}

% Буквы для приложений (исключены Ё, З, Й, О, Ч, Ъ, Ы, Ь — п. 4.11.2)
\def\cyrtoc#1{\ifcase #1\or А\or Б\or В\or Г\or Д\or Е\or Ж\or И\or К\or Л\or М\or Н\or П\or Р\or С\or Т\or У\or Ф\or Х\or Ц\or Ш\or Щ\or Э\or Ю\or Я\else \@ctrerr \fi}

\newcommand\originalappendix{}
\let\originalappendix\appendix
\renewcommand{\appendix}{%
  \originalappendix
  \renewcommand{\thechapter}{\cyrtoc{\value{chapter}}}
  \renewcommand{\appendixname}{ПРИЛОЖЕНИЕ}
  % Заголовок приложения: «ПРИЛОЖЕНИЕ А» по центру, затем название (п. 4.11.2)
  \titleformat{\chapter}{\fontsize{16pt}{18pt}\selectfont\bfseries\centering}{\MakeUppercase{\chaptertitlename}.\ \thechapter}{4pt}{}
  % В приложениях оглавление сокращается до уровня глав
  \addtocontents{toc}{\protect\value{tocdepth}=0}
}

% Переопределяем \mainmatter: убираем сброс счётчика страниц (п. 4.3.1 — сквозная нумерация)
\makeatletter
\renewcommand{\mainmatter}{%
  \cleardoublepage
  \@mainmattertrue
}
\makeatother
